\scanpage{12}
Donnée de $q$ équivaut à la donnée de $\sigma\in\operatorname{Aut}^{\pm1}(\unsure{G|Z},T)$
\[
\left\{\begin{aligned}p(\sigma)&=-1\\ \sigma^2&=1\end{aligned}\right.
\]
\[
\begin{aligned}
q_t(-1)&=\sigma_t\\
\sigma_{tg}&=g\,\sigma_t\,g^{-1}
=g\,\underbrace{\sigma_t\,g^{-1}\,\sigma_t^{-1}}_{g}\,\sigma_t
=g^2\,\sigma_t=\sigma_t\,g^{-2}\\
&\hspace{9em}\boxed{\sigma_{tg}=g^2\sigma_t=\sigma_t g^{-2}}
\end{aligned}
\]

\uline{Corollaire} Si $n$ impair, $t\longmapsto\deleted{\unsure{q}}\sigma_t$ est une bijection de $T$ avec l'ens \deleted{\unsure{des}} $p^{-1}\{-1\}$; donc si on a un \unsure{pol.} à un nb impair de \unsure{cotés|côtés}, l'applic\textsuperscript{n} $t\longmapsto\sigma_t$ qui, à tt sommet du polygone, associe l'unique \unsure{autom.} non orienté qui fixe ce sommet, est \unsure{bijective|bijection}.

\uline{NB} Si $n$ pair, \deleted{\unsure{elle}} ce n'est pas \unsure{bijectif}, \deleted{\unsure{mais}} \deleted{\unsure{pour tt}} ses fibres sont vides ou de cardinal 2\ldots

% slanted
\margin{\tikz[baseline=(side.base)]{\node[inner sep=1pt] (side) {\parbox{.92\marginparwidth}{$\{u,\bar u\}\subset Z$ \unsure{gén.} opposés d'ordre $n=2m$, alors $u^m=\bar u^m=s$, \uline{unique} \uline{\unsure{élémt} d'ordre} 2 de $Z$ appelé symétrie \unsure{Mon} $\sigma_t=\sigma_{t'}$ ssi $t,t'$ sont \unsure{conjugués} mod $\{1,s\}$}}; \draw (side.north east) .. controls +(0,-1em) and +(-.6em,0) .. ([xshift=.45em]side.east) .. controls +(-.7em,0) and +(1em,1.8em) .. (side.south east);}}

\uline{Sommet} \uline{opposé}, symétrie canonique

\uline{Parenthèse}\qquad Produit fibré d'ens

\noindent\begin{minipage}[c]{.46\linewidth}
\diagram{\begin{tikzcd}[ampersand replacement=\&]
E\arrow[r,"p"] \& \unsure{S|B}\\
P\arrow[u,dashed,"q'"]\arrow[r,dashed,"p'"] \& F\arrow[u,"q"']
\end{tikzcd}}
\end{minipage}\hfill
\begin{minipage}[c]{.49\linewidth}
\[
\begin{aligned}
&P\subset E\times F\\
&\hspace{.6em}\rotatebox{90}{$=$}\\
&\{(x,y)\mid px=qy\}
\end{aligned}
\]
\end{minipage}

\unsure{$s$}

Solution d'un Pb universel dans la cat. des ens :

Si $E,F,S$ des groupes, $p,q$ hom de grpe, alors $P$ un groupe (et solution du pb univ. analogue ds cat. des grpes)
\[
\operatorname{Ker}p\sim\operatorname{Ker}p',\quad
\operatorname{Ker}q\sim\operatorname{Ker}q'
\]

Si $p$\inserted{$:E\mathbin{\text{---}}S$} surjectif \unsure{(}bijectif) alors $P\mathbin{\text{---}}F$ surjectif (bijectif).

[Langage des \deleted{\unsure{diagrammes}}\inserted{carrés} cartésiens]

Image inverse d'\unsure{ex.|ext.} de \unsure{groupe|groupes}
% archive sheet number: 11

